\section{Related Work}
\label{sec:related-work}
Prior research relevant to this system falls into three primary domains: data-driven and knowledge-graph approaches to drug-drug interaction (DDI) prediction, multilingual speech interfaces in healthcare, and optical prescription digitization systems.

\subsection{Drug Interaction Databases and Knowledge Graphs}
Early data-driven approaches to identifying drug interactions relied on mining spontaneous adverse event reporting databases alongside pharmacological and side-effect similarity metrics \cite{b2}. While effective at predicting candidate interactions at scale for research, these methods output complex tabular data rather than patient-queryable interfaces. Recent works leverage knowledge graphs (KGs) to represent drugs, side effects, and therapeutic pathways as interconnected entities, utilizing graph neural networks and graph embedding strategies to predict novel drug-drug interactions (DDIs) \cite{b12}, \cite{b15}. Standard clinical repositories such as DrugBank \cite{b1} provide foundational relational datasets for decision support. However, most existing KG frameworks are evaluated on large pharmacological databases aimed at clinicians or researchers \cite{b12}, \cite{b15}. This work differs by anchoring a scoped, verified drug dataset within a hybrid data layer, prioritizing deterministic lookup, explainable graph traversal, and direct patient accessibility over large-scale interaction discovery.

\subsection{Multilingual Speech Recognition and Low-Resource Dialects}
Voice-assisted healthcare tools have gained traction to improve health literacy and bridge communication gaps for disabled and elderly populations \cite{b1}, \cite{b8}, \cite{b9}. Recent advancements in automatic speech recognition (ASR)—including large multilingual transformer models like OpenAI Whisper—have significantly reduced transcription error rates across high-resource languages. However, unwritten regional dialects such as Tulu lack dedicated acoustic training corpora \cite{b10}. To address linguistic barriers in rural healthcare settings, researchers have explored task-oriented voice bots and conversational agents \cite{b5}, \cite{b6}, \cite{b9}, \cite{b10}, \cite{b14}. This work specifically accounts for low-resource constraints by employing a primary multi-locale Google ASR pipeline (`tcy-IN`, `kn-IN`), secondary Whisper ASR fallback, and phonetic Kannada text-to-speech synthesis proxies for Tulu.

\subsection{Optical Prescription Digitization and Label Recognition}
To digitize prescription labels and overcome typing friction, computer vision and Optical Character Recognition (OCR) systems have been widely implemented \cite{b3}, \cite{b4}, \cite{b11}, \cite{b13}, \cite{b16}. Systems utilizing Tesseract OCR \cite{b4} or dual-engine fusion algorithms \cite{b13} enable the extraction of drug names, dosages, and instructions from packaging or cylindrically distorted labels \cite{b11}. Advanced deep learning OCR pipelines combining CNN and RNN architectures further assist in handwritten prescription analysis \cite{b3}, \cite{b16}. However, standalone OCR systems focus strictly on text extraction and do not evaluate extracted drug combinations against structured safety databases.

This work integrates these three lines of research into a unified platform, combining optical prescription parsing with a regional Dravidian voice interface and a deterministic interaction-checking data layer.

